% ============================================
% 外部图片模板 (External Image Templates)
% 适用场景：截图、照片、第三方工具生成的图表
% ============================================

% --------------------------------------------
% 基础模板
% --------------------------------------------
\begin{figure}[H]
  \centering
  \includegraphics[width=0.8\textwidth]{figures/文件名.png}
  \caption{图片说明}
  \label{fig:xxx}
\end{figure}

% --------------------------------------------
% 不同宽度示例
% --------------------------------------------

% 小图（适合简单图标、小截图）
\begin{figure}[H]
  \centering
  \includegraphics[width=0.5\textwidth]{figures/small_image.png}
  \caption{小尺寸图片示例}
  \label{fig:small}
\end{figure}

% 中图（默认推荐）
\begin{figure}[H]
  \centering
  \includegraphics[width=0.7\textwidth]{figures/medium_image.png}
  \caption{中等尺寸图片示例}
  \label{fig:medium}
\end{figure}

% 大图（适合需要看清细节的截图）
\begin{figure}[H]
  \centering
  \includegraphics[width=0.9\textwidth]{figures/large_image.png}
  \caption{大尺寸图片示例}
  \label{fig:large}
\end{figure}

% 全宽（慎用，可能溢出）
\begin{figure}[H]
  \centering
  \includegraphics[width=\textwidth]{figures/full_width.png}
  \caption{全宽图片示例}
  \label{fig:full}
\end{figure}

% --------------------------------------------
% 指定高度（当宽度不好控制时）
% --------------------------------------------
\begin{figure}[H]
  \centering
  \includegraphics[height=6cm]{figures/tall_image.png}
  \caption{按高度缩放的图片}
  \label{fig:by_height}
\end{figure}

% --------------------------------------------
% 并排双图
% --------------------------------------------
\begin{figure}[H]
  \centering
  \begin{minipage}{0.45\textwidth}
    \centering
    \includegraphics[width=\textwidth]{figures/image_a.png}
    \subcaption{图片A说明}
    \label{fig:side_a}
  \end{minipage}
  \hfill
  \begin{minipage}{0.45\textwidth}
    \centering
    \includegraphics[width=\textwidth]{figures/image_b.png}
    \subcaption{图片B说明}
    \label{fig:side_b}
  \end{minipage}
  \caption{并排图片的整体说明}
  \label{fig:side_by_side}
\end{figure}

% --------------------------------------------
% 支持的图片格式
% --------------------------------------------
% .png  - 推荐，适合截图、有透明背景的图
% .jpg  - 适合照片
% .pdf  - 适合矢量图，缩放不失真